\documentclass[11pt,a4paper]{article}

% ================ Encodage et langue ================
\usepackage[utf8]{inputenc}
\usepackage[T1]{fontenc}
\usepackage[french]{babel}
\usepackage{lmodern}
\usepackage{textcomp}

% ================ Mise en page ================
\usepackage[margin=2.2cm]{geometry}
\usepackage{parskip}
\setlength{\parindent}{0pt}
\setlength{\parskip}{0.6em}

% ================ Couleurs et liens ================
\usepackage[table]{xcolor}
\definecolor{primary}{HTML}{0B5394}
\definecolor{accent}{HTML}{1F77B4}
\definecolor{light}{HTML}{F2F2F2}
\definecolor{codebg}{HTML}{F6F8FA}
\definecolor{rulec}{HTML}{D0D7DE}

\usepackage[colorlinks=true, linkcolor=primary, urlcolor=primary, citecolor=primary]{hyperref}

% ================ Tableaux ================
\usepackage{longtable, booktabs, array, tabularx}
\renewcommand{\arraystretch}{1.25}
\arrayrulecolor{rulec}

% ================ Listes ================
\usepackage{enumitem}
\setlist{itemsep=2pt, topsep=4pt}

% ================ Code et citations ================
\usepackage{fancyvrb}
\usepackage{framed}
\definecolor{shadecolor}{named}{light}

% Encadrés pour les blocs > (quotes)
\usepackage{mdframed}
\newmdenv[
  linecolor=primary,
  linewidth=2pt,
  topline=false, bottomline=false, rightline=false,
  innerleftmargin=10pt, innerrightmargin=8pt,
  innertopmargin=6pt, innerbottommargin=6pt,
  backgroundcolor=codebg,
  skipabove=8pt, skipbelow=8pt
]{quotebox}
\renewenvironment{quote}{\begin{quotebox}}{\end{quotebox}}

% ================ Titres ================
\usepackage{titlesec}
\titleformat{\section}{\Large\bfseries\color{primary}}{\thesection.}{0.6em}{}
\titleformat{\subsection}{\large\bfseries\color{accent}}{\thesubsection}{0.6em}{}
\titleformat{\subsubsection}{\normalsize\bfseries}{\thesubsubsection}{0.6em}{}

% ================ En-têtes et pieds de page ================
\usepackage{fancyhdr}
\pagestyle{fancy}
\fancyhf{}
\fancyhead[L]{\small\color{gray}ESPRIT School of Business --- M1 BA --- Projet Intégré}
\fancyhead[R]{\small\color{gray}\nouppercase{\leftmark}}
\fancyfoot[C]{\small\thepage}
\renewcommand{\headrulewidth}{0.3pt}
\renewcommand{\footrulewidth}{0pt}

% ================ Pandoc helpers ================
\providecommand{\tightlist}{\setlength{\itemsep}{0pt}\setlength{\parskip}{0pt}}
\providecommand{\passthrough}[1]{#1}

\usepackage{calc}
\providecommand{\real}[1]{#1}


% ================ Métadonnées ================
\title{\color{primary}\textbf{Guide Étudiant --- Démarche et Outils par
Étape}}
\author{Aymen Ben Brik}
\date{2026}

\begin{document}

\begin{titlepage}
\centering
\vspace*{1.5cm}

{\Large\bfseries ESPRIT School of Business}\\[0.3cm]
{\large Master 1 --- Business Analytics (M1 BA)}\\[2cm]

{\color{gray}\large Projet Intégré}\\[0.3cm]
{\color{gray}\large Analyse et Prédiction du Churn Client}\\[2.2cm]

\rule{0.85\linewidth}{1.2pt}\\[0.6cm]
{\Huge\bfseries\color{primary} Guide Étudiant --- Démarche et Outils par
Étape\par}
\rule{0.85\linewidth}{1.2pt}\\[2cm]

{\large Tuteur du projet}\\[0.2cm]
{\Large\bfseries Aymen Ben Brik}\\[1cm]
{\large 2026}

\vfill
{\color{gray}\small Document à destination des étudiants}
\end{titlepage}

\tableofcontents
\newpage

\section{Guide Étudiant --- Démarche et Outils par
Étape}\label{guide-uxe9tudiant-duxe9marche-et-outils-par-uxe9tape}

Ce guide vous accompagne \textbf{étape par étape} dans la réalisation du
projet. Pour chaque phase, vous trouverez :

\begin{itemize}
\tightlist
\item
  Les \textbf{objectifs} à atteindre.
\item
  Les \textbf{livrables} intermédiaires attendus.
\item
  Une \textbf{proposition d'outils} parmi lesquels choisir (vous restez
  libres, mais justifiez vos choix dans le rapport).
\end{itemize}

\begin{quote}
\textbf{Conseil général} : ne cherchez pas à utiliser les outils les
plus complexes. Privilégiez ceux que vous maîtrisez ou que vous pouvez
prendre en main rapidement. Un outil simple bien utilisé vaut mieux
qu'un outil sophistiqué mal utilisé.
\end{quote}

\begin{center}\rule{0.5\linewidth}{0.5pt}\end{center}

\subsection{Étape 0 --- Mise en place de
l'environnement}\label{uxe9tape-0-mise-en-place-de-lenvironnement}

\subsubsection{Objectifs}\label{objectifs}

\begin{itemize}
\tightlist
\item
  Préparer un environnement de travail reproductible.
\item
  Mettre en place le dépôt GitHub et la collaboration.
\end{itemize}

\subsubsection{À faire}\label{uxe0-faire}

\begin{itemize}
\tightlist
\item
  Créer un dépôt \textbf{GitHub} privé pour la phase de développement,
  public à la livraison (ou public dès le départ si vous excluez bien
  les données).
\item
  Inviter l'encadrant en tant que collaborateur.
\item
  Créer un environnement virtuel Python (\texttt{venv} ou
  \texttt{conda}).
\item
  Rédiger un \texttt{requirements.txt} ou \texttt{environment.yml}.
\item
  Rédiger un premier \texttt{README.md} avec : description, structure du
  dépôt, instructions d'installation.
\item
  Configurer un \texttt{.gitignore} qui exclut \textbf{toutes les
  données} (\texttt{*.csv}, \texttt{*.xlsx}, \texttt{*.pkl}, dossier
  \texttt{data/}, etc.) et les fichiers d'environnement (\texttt{venv/},
  \texttt{\_\_pycache\_\_/}, \texttt{.env}).
\end{itemize}

\subsubsection{Outils proposés}\label{outils-proposuxe9s}

{\def\LTcaptype{none} % do not increment counter
\begin{longtable}[]{@{}ll@{}}
\toprule\noalign{}
Catégorie & Options \\
\midrule\noalign{}
\endhead
\bottomrule\noalign{}
\endlastfoot
Versioning & \textbf{Git + GitHub} (obligatoire) \\
IDE & VS Code, PyCharm, JetBrains DataSpell \\
Environnement Python & \texttt{venv}, \texttt{conda}, \texttt{mamba},
\texttt{uv}, \texttt{poetry} \\
Notebooks & Jupyter Lab, Jupyter Notebook, VS Code Notebooks \\
Gestion de projet & GitHub Projects, Trello, Notion \\
\end{longtable}
}

\begin{center}\rule{0.5\linewidth}{0.5pt}\end{center}

\subsection{Étape 1 --- Exploration des données
(EDA)}\label{uxe9tape-1-exploration-des-donnuxe9es-eda}

\subsubsection{Objectifs}\label{objectifs-1}

\begin{itemize}
\tightlist
\item
  Comprendre la structure, la volumétrie et la qualité des données.
\item
  Identifier les variables intéressantes pour le churn.
\item
  Documenter les problèmes de qualité à traiter.
\end{itemize}

\subsubsection{À faire}\label{uxe0-faire-1}

\begin{itemize}
\tightlist
\item
  Charger les fichiers et vérifier leur intégrité.
\item
  Calculer : nombre de lignes, types des colonnes, taux de valeurs
  manquantes, distributions.
\item
  Croiser les données avec les dimensions (vérifier que les clés
  correspondent).
\item
  Visualiser : distributions par variable, corrélations, taux de churn
  par segment.
\item
  Produire un \textbf{rapport d'exploration} (notebook) avec vos
  premières observations.
\end{itemize}

\subsubsection{Outils proposés}\label{outils-proposuxe9s-1}

{\def\LTcaptype{none} % do not increment counter
\begin{longtable}[]{@{}
  >{\raggedright\arraybackslash}p{(\linewidth - 2\tabcolsep) * \real{0.5000}}
  >{\raggedright\arraybackslash}p{(\linewidth - 2\tabcolsep) * \real{0.5000}}@{}}
\toprule\noalign{}
\begin{minipage}[b]{\linewidth}\raggedright
Catégorie
\end{minipage} & \begin{minipage}[b]{\linewidth}\raggedright
Options
\end{minipage} \\
\midrule\noalign{}
\endhead
\bottomrule\noalign{}
\endlastfoot
Manipulation & \textbf{pandas}, \textbf{polars} (plus rapide sur grosse
volumétrie), DuckDB \\
Profilage automatique & \texttt{ydata-profiling} (ex-pandas-profiling),
\texttt{sweetviz}, \texttt{dataprep} \\
Visualisation & \textbf{matplotlib}, \textbf{seaborn}, \textbf{plotly},
\textbf{altair} \\
\end{longtable}
}

\begin{center}\rule{0.5\linewidth}{0.5pt}\end{center}

\subsection{Étape 2 --- ETL (Extract -- Transform --
Load)}\label{uxe9tape-2-etl-extract-transform-load}

\subsubsection{Objectifs}\label{objectifs-2}

\begin{itemize}
\tightlist
\item
  Construire un pipeline reproductible qui transforme les données brutes
  en données analytiques propres.
\item
  Charger ces données dans une base analytique exploitable par la BI et
  le ML.
\end{itemize}

\subsubsection{À faire}\label{uxe0-faire-2}

\begin{enumerate}
\def\labelenumi{\arabic{enumi}.}
\tightlist
\item
  \textbf{Extract} : lecture des CSV et fichiers Excel.
\item
  \textbf{Transform} :

  \begin{itemize}
  \tightlist
  \item
    Conversion des dates (\texttt{YYYYMMDD} → \texttt{datetime}).
  \item
    Traitement des valeurs \texttt{NULL} textuelles.
  \item
    Jointures avec les tables de dimensions.
  \item
    Calcul de variables dérivées (âge, ancienneté, etc.).
  \item
    Déduplication et contrôle de cohérence.
  \end{itemize}
\item
  \textbf{Load} : chargement dans le Data Warehouse cible.
\end{enumerate}

\subsubsection{Outils proposés}\label{outils-proposuxe9s-2}

Trois familles d'outils sont acceptées pour l'intégration des données.
Choisissez celle qui correspond le mieux aux compétences de l'équipe et
aux contraintes de votre poste de travail.

{\def\LTcaptype{none} % do not increment counter
\begin{longtable}[]{@{}
  >{\raggedright\arraybackslash}p{(\linewidth - 6\tabcolsep) * \real{0.2500}}
  >{\raggedright\arraybackslash}p{(\linewidth - 6\tabcolsep) * \real{0.2500}}
  >{\raggedright\arraybackslash}p{(\linewidth - 6\tabcolsep) * \real{0.2500}}
  >{\raggedright\arraybackslash}p{(\linewidth - 6\tabcolsep) * \real{0.2500}}@{}}
\toprule\noalign{}
\begin{minipage}[b]{\linewidth}\raggedright
Outil
\end{minipage} & \begin{minipage}[b]{\linewidth}\raggedright
Type
\end{minipage} & \begin{minipage}[b]{\linewidth}\raggedright
Forces
\end{minipage} & \begin{minipage}[b]{\linewidth}\raggedright
À considérer
\end{minipage} \\
\midrule\noalign{}
\endhead
\bottomrule\noalign{}
\endlastfoot
\textbf{Talend Open Studio} & ETL graphique open source & Interface
visuelle (drag \& drop), nombreux connecteurs, génère du code Java,
bonne traçabilité du flux & Installation lourde (JDK requis), démarrage
plus long, version Open Studio non maintenue récemment \\
\textbf{SSIS} (SQL Server Integration Services) & ETL graphique
Microsoft & Très bien intégré avec SQL Server et Power BI, mature en
entreprise, performant sur grosse volumétrie & Nécessite Windows +
Visual Studio + SQL Server Data Tools, écosystème Microsoft requis \\
\textbf{Python} (pandas / polars + SQLAlchemy / DuckDB) & ETL
programmatique & Flexible, reproductible, versionnable facilement,
transition naturelle vers le ML & Pas d'interface graphique, demande de
la rigueur dans la structuration du code \\
\end{longtable}
}

\paragraph{Comment choisir ?}\label{comment-choisir}

\begin{itemize}
\tightlist
\item
  \textbf{Vous visez l'écosystème Microsoft (Power BI + SQL Server)} →
  \textbf{SSIS} est cohérent de bout en bout.
\item
  \textbf{Vous voulez une approche visuelle indépendante de la base
  cible} → \textbf{Talend Open Studio}.
\item
  \textbf{Vous voulez la solution la plus rapide à mettre en place et
  compatible avec le ML qui suit} → \textbf{Python}.
\end{itemize}

\paragraph{Outils complémentaires}\label{outils-compluxe9mentaires}

{\def\LTcaptype{none} % do not increment counter
\begin{longtable}[]{@{}
  >{\raggedright\arraybackslash}p{(\linewidth - 2\tabcolsep) * \real{0.5000}}
  >{\raggedright\arraybackslash}p{(\linewidth - 2\tabcolsep) * \real{0.5000}}@{}}
\toprule\noalign{}
\begin{minipage}[b]{\linewidth}\raggedright
Catégorie
\end{minipage} & \begin{minipage}[b]{\linewidth}\raggedright
Options
\end{minipage} \\
\midrule\noalign{}
\endhead
\bottomrule\noalign{}
\endlastfoot
Transformation SQL avancée & \textbf{dbt}, \textbf{DuckDB}, SQL natif \\
Orchestrateur (optionnel, \emph{overkill} sur 1 mois) & Apache Airflow,
Prefect, Dagster \\
Alternative no-code & Pentaho Kettle, KNIME \\
\end{longtable}
}

\begin{quote}
\textbf{Important} : quel que soit l'outil retenu, votre pipeline doit
être \textbf{reproductible} (un tiers doit pouvoir le ré-exécuter) et
\textbf{documenté} par un schéma clair du flux de données.
\end{quote}

\begin{center}\rule{0.5\linewidth}{0.5pt}\end{center}

\subsection{Étape 3 --- Data Warehouse et modélisation dimensionnelle
(BI)}\label{uxe9tape-3-data-warehouse-et-moduxe9lisation-dimensionnelle-bi}

\subsubsection{Objectifs}\label{objectifs-3}

\begin{itemize}
\tightlist
\item
  Stocker les données préparées dans une base interrogeable.
\item
  Modéliser en étoile : une table de faits + plusieurs tables de
  dimensions.
\end{itemize}

\subsubsection{À faire}\label{uxe0-faire-3}

\begin{itemize}
\tightlist
\item
  Choisir la base de données cible.
\item
  Modéliser la table de faits : grain = (client, compte, mois) ou
  (client, compte) selon votre choix.
\item
  Définir les dimensions : \texttt{dim\_client}, \texttt{dim\_compte},
  \texttt{dim\_produit}, \texttt{dim\_temps}, \texttt{dim\_industrie},
  \texttt{dim\_devise}, \texttt{dim\_motif\_cloture}, etc.
\item
  Créer les scripts SQL de création des tables.
\item
  Charger les données depuis l'ETL.
\item
  Rédiger un schéma visuel (diagramme en étoile).
\end{itemize}

\subsubsection{Outils proposés}\label{outils-proposuxe9s-3}

{\def\LTcaptype{none} % do not increment counter
\begin{longtable}[]{@{}
  >{\raggedright\arraybackslash}p{(\linewidth - 2\tabcolsep) * \real{0.5000}}
  >{\raggedright\arraybackslash}p{(\linewidth - 2\tabcolsep) * \real{0.5000}}@{}}
\toprule\noalign{}
\begin{minipage}[b]{\linewidth}\raggedright
Catégorie
\end{minipage} & \begin{minipage}[b]{\linewidth}\raggedright
Options
\end{minipage} \\
\midrule\noalign{}
\endhead
\bottomrule\noalign{}
\endlastfoot
Base de données analytique légère & \textbf{DuckDB} (excellent pour ce
projet : zéro configuration), \textbf{SQLite} \\
Base de données SQL classique & \textbf{PostgreSQL}, \textbf{MySQL},
\textbf{MariaDB} \\
Base de données analytique avancée & \textbf{ClickHouse}, \textbf{MS SQL
Server} (utile si vous visez Power BI sur site) \\
Modélisation dimensionnelle & \textbf{dbt} (transformation + tests), SQL
natif \\
Visualisation du schéma & dbdiagram.io, draw.io, Lucidchart, Mermaid \\
\end{longtable}
}

\begin{quote}
\textbf{Recommandation} : \textbf{DuckDB} ou \textbf{PostgreSQL} sont
les meilleurs choix pour ce projet. DuckDB pour la simplicité,
PostgreSQL si vous voulez une expérience plus proche de la production.
\end{quote}

\begin{center}\rule{0.5\linewidth}{0.5pt}\end{center}

\subsection{Étape 4 --- Tableaux de bord Power
BI}\label{uxe9tape-4-tableaux-de-bord-power-bi}

\subsubsection{Objectifs}\label{objectifs-4}

\begin{itemize}
\tightlist
\item
  Restituer les KPIs métier de manière interactive.
\item
  Permettre une exploration par profil utilisateur (direction,
  conseillers, marketing).
\end{itemize}

\subsubsection{À faire}\label{uxe0-faire-4}

\begin{itemize}
\tightlist
\item
  Se connecter à la source de données (DWH).
\item
  Définir les mesures DAX nécessaires (taux de churn, ancienneté
  moyenne, etc.).
\item
  Construire au minimum \textbf{3 pages} :

  \begin{enumerate}
  \def\labelenumi{\arabic{enumi}.}
  \tightlist
  \item
    \textbf{Vue d'ensemble} : KPIs clés, taux de churn global,
    évolution.
  \item
    \textbf{Analyse du churn} : décomposition par segment (âge, secteur,
    ancienneté, produit, devise).
  \item
    \textbf{Segmentation client} : profils, valeur client, comptes à
    risque.
  \end{enumerate}
\item
  Soigner la mise en page : filtres, slicers, navigation entre pages,
  infobulles.
\item
  Publier le rapport (Power BI Service ou fichier \texttt{.pbix}).
\end{itemize}

\subsubsection{Outils proposés}\label{outils-proposuxe9s-4}

{\def\LTcaptype{none} % do not increment counter
\begin{longtable}[]{@{}
  >{\raggedright\arraybackslash}p{(\linewidth - 2\tabcolsep) * \real{0.5000}}
  >{\raggedright\arraybackslash}p{(\linewidth - 2\tabcolsep) * \real{0.5000}}@{}}
\toprule\noalign{}
\begin{minipage}[b]{\linewidth}\raggedright
Catégorie
\end{minipage} & \begin{minipage}[b]{\linewidth}\raggedright
Options
\end{minipage} \\
\midrule\noalign{}
\endhead
\bottomrule\noalign{}
\endlastfoot
Outil principal & \textbf{Power BI Desktop} (obligatoire --- l'énoncé
l'impose) \\
Publication & Power BI Service (compte gratuit pour la publication,
payant pour le partage) \\
Alternatives complémentaires & Tableau, Metabase, Apache Superset,
Looker Studio (non --- Power BI exigé) \\
\end{longtable}
}

\begin{quote}
\textbf{Conseil} : la qualité du design des dashboards compte.
Inspirez-vous des galeries Power BI officielles. Ne surchargez pas vos
pages.
\end{quote}

\begin{center}\rule{0.5\linewidth}{0.5pt}\end{center}

\subsection{Étape 5 --- Machine
Learning}\label{uxe9tape-5-machine-learning}

\subsubsection{Objectifs}\label{objectifs-5}

\begin{itemize}
\tightlist
\item
  Construire un modèle qui prédit la probabilité de churn d'un client.
\item
  Choisir le meilleur modèle de manière justifiée.
\end{itemize}

\subsubsection{À faire}\label{uxe0-faire-5}

\begin{enumerate}
\def\labelenumi{\arabic{enumi}.}
\tightlist
\item
  \textbf{Préparation} : sélectionner les features, encoder les
  variables catégorielles, gérer le déséquilibre (poids des classes,
  SMOTE, etc.), split train/test stratifié.
\item
  \textbf{Modélisation} : entraîner \textbf{au moins 3 modèles}
  différents (par exemple : régression logistique en baseline, Random
  Forest, XGBoost).
\item
  \textbf{Évaluation} : utiliser les métriques adaptées au déséquilibre
  (precision, recall, F1, \textbf{ROC-AUC}, \textbf{PR-AUC}). Pas
  seulement l'accuracy.
\item
  \textbf{Tuning} : optimiser les hyperparamètres du meilleur candidat.
\item
  \textbf{Interprétation} : importance des features + valeurs SHAP.
\item
  \textbf{Sauvegarde} : sérialiser le modèle final.
\end{enumerate}

\subsubsection{Outils proposés}\label{outils-proposuxe9s-5}

{\def\LTcaptype{none} % do not increment counter
\begin{longtable}[]{@{}
  >{\raggedright\arraybackslash}p{(\linewidth - 2\tabcolsep) * \real{0.5000}}
  >{\raggedright\arraybackslash}p{(\linewidth - 2\tabcolsep) * \real{0.5000}}@{}}
\toprule\noalign{}
\begin{minipage}[b]{\linewidth}\raggedright
Catégorie
\end{minipage} & \begin{minipage}[b]{\linewidth}\raggedright
Options
\end{minipage} \\
\midrule\noalign{}
\endhead
\bottomrule\noalign{}
\endlastfoot
Manipulation & \textbf{pandas}, \textbf{numpy} \\
Modèles classiques & \textbf{scikit-learn} \\
Modèles boostés & \textbf{XGBoost}, \textbf{LightGBM},
\textbf{CatBoost} \\
Déséquilibre de classes & \texttt{imbalanced-learn} (SMOTE),
\texttt{scale\_pos\_weight} \\
Hyperparamètres & \texttt{GridSearchCV}, \texttt{RandomizedSearchCV},
\textbf{Optuna} \\
Interprétabilité & \textbf{SHAP}, \texttt{eli5}, importance native des
arbres \\
Sauvegarde & \textbf{joblib}, \texttt{pickle}, formats ONNX ou MLflow \\
Suivi d'expériences (optionnel) & \textbf{MLflow}, Weights \& Biases \\
\end{longtable}
}

\begin{center}\rule{0.5\linewidth}{0.5pt}\end{center}

\subsection{Étape 6 --- Interface web et
déploiement}\label{uxe9tape-6-interface-web-et-duxe9ploiement}

\subsubsection{Objectifs}\label{objectifs-6}

\begin{itemize}
\tightlist
\item
  Exposer le modèle ML et les analyses à des utilisateurs
  non-techniques.
\item
  Déployer en ligne pour une démonstration accessible.
\end{itemize}

\subsubsection{À faire}\label{uxe0-faire-6}

\begin{itemize}
\tightlist
\item
  Concevoir une interface simple avec au moins :

  \begin{itemize}
  \tightlist
  \item
    Un \textbf{formulaire de prédiction individuelle} (saisir un profil
    client → recevoir un score de churn).
  \item
    Une \textbf{liste des clients actifs à risque} issue d'une
    prédiction en lot.
  \item
    Des \textbf{KPIs de synthèse}.
  \end{itemize}
\item
  Brancher le modèle sauvegardé à l'application.
\item
  Déployer sur une plateforme cloud accessible publiquement (URL fournie
  dans le README).
\end{itemize}

\subsubsection{Outils proposés --- Front + back unifié (recommandé pour
ce
projet)}\label{outils-proposuxe9s-front-back-unifiuxe9-recommanduxe9-pour-ce-projet}

{\def\LTcaptype{none} % do not increment counter
\begin{longtable}[]{@{}
  >{\raggedright\arraybackslash}p{(\linewidth - 4\tabcolsep) * \real{0.3333}}
  >{\raggedright\arraybackslash}p{(\linewidth - 4\tabcolsep) * \real{0.3333}}
  >{\raggedright\arraybackslash}p{(\linewidth - 4\tabcolsep) * \real{0.3333}}@{}}
\toprule\noalign{}
\begin{minipage}[b]{\linewidth}\raggedright
Outil
\end{minipage} & \begin{minipage}[b]{\linewidth}\raggedright
Forces
\end{minipage} & \begin{minipage}[b]{\linewidth}\raggedright
Limites
\end{minipage} \\
\midrule\noalign{}
\endhead
\bottomrule\noalign{}
\endlastfoot
\textbf{Streamlit} & Très rapide à développer en Python pur, idéal pour
démo data & Personnalisation graphique limitée \\
\textbf{Gradio} & Très simple pour exposer un modèle ML & Moins
polyvalent qu'un dashboard complet \\
\textbf{Dash (Plotly)} & Très bonnes visualisations, plus contrôlable &
Courbe d'apprentissage un peu plus longue \\
\end{longtable}
}

\subsubsection{Outils proposés --- Front et back séparés (si vous voulez
aller plus
loin)}\label{outils-proposuxe9s-front-et-back-suxe9paruxe9s-si-vous-voulez-aller-plus-loin}

{\def\LTcaptype{none} % do not increment counter
\begin{longtable}[]{@{}ll@{}}
\toprule\noalign{}
Couche & Options \\
\midrule\noalign{}
\endhead
\bottomrule\noalign{}
\endlastfoot
API back-end & \textbf{FastAPI} (recommandé), Flask, Django REST \\
Front-end & React, Vue, Angular, Next.js, simple HTML/Bootstrap \\
\end{longtable}
}

\subsubsection{Outils proposés --- Hébergement /
Déploiement}\label{outils-proposuxe9s-huxe9bergement-duxe9ploiement}

{\def\LTcaptype{none} % do not increment counter
\begin{longtable}[]{@{}
  >{\raggedright\arraybackslash}p{(\linewidth - 2\tabcolsep) * \real{0.5000}}
  >{\raggedright\arraybackslash}p{(\linewidth - 2\tabcolsep) * \real{0.5000}}@{}}
\toprule\noalign{}
\begin{minipage}[b]{\linewidth}\raggedright
Option
\end{minipage} & \begin{minipage}[b]{\linewidth}\raggedright
Idéal pour
\end{minipage} \\
\midrule\noalign{}
\endhead
\bottomrule\noalign{}
\endlastfoot
\textbf{Streamlit Community Cloud} & App Streamlit (gratuit, ultra
simple) \\
\textbf{Hugging Face Spaces} & App Gradio / Streamlit (gratuit) \\
\textbf{Render} & App Python complète (free tier) \\
\textbf{Railway} & App Python + DB (free tier limité) \\
\textbf{Vercel / Netlify} & Front statique ou Next.js \\
\textbf{Docker + un cloud (Azure, GCP, AWS)} & Si vous maîtrisez
Docker \\
\end{longtable}
}

\begin{quote}
\textbf{Recommandation} : \textbf{Streamlit + Streamlit Cloud} ou
\textbf{Gradio + Hugging Face Spaces} pour un déploiement rapide et
gratuit. Visez la simplicité.
\end{quote}

\begin{center}\rule{0.5\linewidth}{0.5pt}\end{center}

\subsection{Étape 7 --- Rapport et
présentation}\label{uxe9tape-7-rapport-et-pruxe9sentation}

\subsubsection{Objectifs}\label{objectifs-7}

\begin{itemize}
\tightlist
\item
  Documenter votre démarche, vos choix et vos résultats.
\item
  Convaincre le jury en 20 minutes que votre solution répond au besoin.
\end{itemize}

\subsubsection{Structure recommandée du rapport (25--40
pages)}\label{structure-recommanduxe9e-du-rapport-2540-pages}

\begin{enumerate}
\def\labelenumi{\arabic{enumi}.}
\tightlist
\item
  \textbf{Page de garde et résumé exécutif} (1 page).
\item
  \textbf{Introduction et contexte} (2 pages).
\item
  \textbf{Description et exploration des données} (3--5 pages).
\item
  \textbf{Architecture de la solution} (2--3 pages, avec schémas).
\item
  \textbf{ETL et modélisation dimensionnelle} (3--5 pages).
\item
  \textbf{Analyses BI et dashboards Power BI} (3--5 pages, captures
  d'écran).
\item
  \textbf{Modélisation ML} (5--8 pages, métriques, comparaisons,
  interprétation).
\item
  \textbf{Interface web et déploiement} (2--3 pages, captures, lien
  démo).
\item
  \textbf{Limites et perspectives} (1--2 pages).
\item
  \textbf{Conclusion} (1 page).
\item
  \textbf{Annexes} : références, code clés, glossaire.
\end{enumerate}

\subsubsection{Outils proposés}\label{outils-proposuxe9s-6}

{\def\LTcaptype{none} % do not increment counter
\begin{longtable}[]{@{}
  >{\raggedright\arraybackslash}p{(\linewidth - 2\tabcolsep) * \real{0.5000}}
  >{\raggedright\arraybackslash}p{(\linewidth - 2\tabcolsep) * \real{0.5000}}@{}}
\toprule\noalign{}
\begin{minipage}[b]{\linewidth}\raggedright
Catégorie
\end{minipage} & \begin{minipage}[b]{\linewidth}\raggedright
Options
\end{minipage} \\
\midrule\noalign{}
\endhead
\bottomrule\noalign{}
\endlastfoot
Rédaction du rapport & \textbf{LaTeX} (Overleaf), \textbf{Microsoft
Word}, \textbf{Google Docs}, \textbf{Typst}, Markdown + Pandoc \\
Diagrammes & \textbf{draw.io}, \textbf{Lucidchart}, \textbf{Mermaid},
\textbf{Excalidraw}, dbdiagram.io \\
Captures d'écran & ShareX, Snipping Tool, macOS Screenshot \\
Slides & \textbf{PowerPoint}, \textbf{Google Slides}, \textbf{Canva},
\textbf{Beamer} (LaTeX), Marp (Markdown) \\
\end{longtable}
}

\subsubsection{Conseils pour la présentation
orale}\label{conseils-pour-la-pruxe9sentation-orale}

\begin{itemize}
\tightlist
\item
  \textbf{Temps} : 20 minutes de présentation + 10 minutes de questions.
  Ne dépassez pas.
\item
  \textbf{Rythme} : \textasciitilde1 slide par minute.
\item
  \textbf{Démo} : prévoyez \textbf{au moins 5 minutes} de démonstration
  live de votre application web. Préparez un plan B (vidéo de secours)
  en cas de problème de connexion.
\item
  \textbf{Répartition} : chaque membre de l'équipe doit prendre la
  parole.
\item
  \textbf{Public} : présentez à un public à la fois technique et métier
  --- expliquez les choix techniques de manière accessible.
\end{itemize}

\begin{center}\rule{0.5\linewidth}{0.5pt}\end{center}

\subsection{Récapitulatif --- Stack recommandée (option simple et
rapide)}\label{ruxe9capitulatif-stack-recommanduxe9e-option-simple-et-rapide}

Si vous voulez une stack qui ``fonctionne'', voici une combinaison
cohérente :

{\def\LTcaptype{none} % do not increment counter
\begin{longtable}[]{@{}ll@{}}
\toprule\noalign{}
Composante & Choix recommandé \\
\midrule\noalign{}
\endhead
\bottomrule\noalign{}
\endlastfoot
Langage & Python 3.10+ \\
EDA / ETL & pandas + Jupyter \\
Data Warehouse & DuckDB ou PostgreSQL \\
Dashboards & Power BI Desktop \\
ML & scikit-learn + XGBoost \\
Interface web & Streamlit \\
Déploiement & Streamlit Community Cloud \\
Versioning & Git + GitHub \\
Rapport & Word ou LaTeX (Overleaf) \\
Slides & PowerPoint ou Google Slides \\
\end{longtable}
}

Cette stack permet de couvrir tout le périmètre en restant dans le temps
imparti. Vous êtes libres d'en choisir une autre, \textbf{mais justifiez
vos choix dans le rapport}.

\end{document}
